\section{Results}
\label{sec:results}

\subsection{What the radio image knows}
\begin{itemize}
  \item Pooled CNN $0.487 \pm 0.033$ (ensemble 0.528) vs total flux 0.432,
        mass 0.314, five scalars 0.455.
  \item Within mass terciles: CNN 0.33/0.32/0.22 while scalars fall to
        $-0.28$ in the high-mass bin; partial $\rho = 0.475$ at fixed mass.
  \item Morphology alone (flux-normalised maps): 0.418; but shape and
        brightness are not separable (shape-only model 59\% predictable
        from the flux it never saw).
  \item Caveat: scalars explain 72\% of the CNN output.
  \item Minkowski functionals as an interpretable baseline: XGBoost on 48
        topology numbers reaches 0.454; fragmentation rises with merger age.
\end{itemize}
\begin{figure*}
  \centering
  \includegraphics[width=\textwidth]{minkowski.png}
  \caption{Minkowski functionals of the simulated radio maps and their
  relation to merger time.
  \stale{built on weight-centroid-centred maps before the density cut and
  the protocol fix; rerun \texttt{minkowski\_functionals.py} on
  \texttt{dataset\_nh4\_128.h5}.}}
  \label{fig:mf}
\end{figure*}
\figplaceholder{Bar chart: OOF $R^2$ of the scalar baselines, morphology
only, CNN, ensemble (\texttt{analyze\_mass\_control.py},
\texttt{compare\_shape.py}).}{The image against scalar baselines.}
\figplaceholder{True vs predicted TSC, with residuals by age bin
(\texttt{plot\_tsc\_predictions.py} on saved OOF predictions).}
{Predicted against true merger time.}

\subsection{X-ray, and radio with X-ray}
\begin{itemize}
  \item The stretch bug and fix: X-ray 0.320 $\to$ 0.511 ($p = 0.0003$).
  \item Joint model $0.537 \pm 0.016$: beats radio by $+0.050$ and X-ray by
        $+0.026$; gains on older mergers (RMSE $>2$\,Gyr 1.25 $\to$ 1.13).
\end{itemize}

\subsection{Under observational degradation}
\label{sec:flux}
\begin{itemize}
  \item Anchored vs simulation-power flux: 0.338 $\to$ 0.412; image beyond
        mass $+0.079 \to +0.151$; $\rho({\rm pred}, M)$ $-0.89 \to -0.61$.
  \item The simulation's $P_{150}$--\mfive: slope $3.30\pm0.33$ (obs.\ 3.55),
        scatter 1.65\,dex (obs.\ 0.35); scatter tracks merger state with
        the observed sign.
  \item Realistic X-ray and joint (after all background fixes): X-ray 0.339,
        joint 0.438. \todo{5 seeds; snapshot 91.}
\end{itemize}
\figplaceholder{Simulated $P_{150}$--\mfive\ relation coloured by TSC,
with the Cuciti et al.\ (2023) fit (\texttt{summarize\_b23.py} data).}
{The simulation's own mass--power relation.}

\subsection{Application to LoVoCCS}
\begin{itemize}
  \item Predictions for 25 (radio), 15 (X-ray), 14 (joint) clusters; all
        within the label range.
  \item Radio and X-ray agree: $\rho = 0.62$ ($p = 0.019$), 0.67 at fixed
        $z$; at fixed weak-lensing mass 0.38 ($n = 10$, n.s.).
  \item A399/A401 youngest -- but mass-driven; do not present as
        validation.
  \item Weak-lensing peak check: only 9 overlapping clusters,
        inconclusive.
\end{itemize}
\figplaceholder{Radio vs X-ray predictions for the 14 clusters with both,
coloured by weak-lensing mass (\texttt{compare\_real\_preds.py}).}
{Two independent observables on real clusters.}
